\documentclass[10pt,a4paper]{amsart}
\usepackage{amsmath,amssymb,amsthm,graphicx}
% For commutative diagrams we use Xy-pic
\usepackage[all]{xy}

% We use 2cell for 2-commutative diagrams.
\xyoption{2cell}
% Theorem environments.
%
\theoremstyle{plain}
\newtheorem{theorem}[subsection]{Theorem}
\newtheorem{proposition}[subsection]{Proposition}
\newtheorem{lemma}[subsection]{Lemma}

\theoremstyle{definition}
\newtheorem{definition}[subsection]{Definition}
\newtheorem{example}[subsection]{Example}
\newtheorem{exercise}[subsection]{Exercise}
\newtheorem{situation}[subsection]{Situation}

\theoremstyle{remark}
\newtheorem{remark}[subsection]{Remark}
\newtheorem{remarks}[subsection]{Remarks}

\numberwithin{equation}{subsection}

% Macros
%
\def\lim{\mathop{\mathrm{lim}}\nolimits}
\def\colim{\mathop{\mathrm{colim}}\nolimits}
\def\Spec{\mathop{\mathrm{Spec}}}
\def\Hom{\mathop{\mathrm{Hom}}\nolimits}
\def\Ext{\mathop{\mathrm{Ext}}\nolimits}
\def\SheafHom{\mathop{\mathcal{H}\!\mathit{om}}\nolimits}
\def\SheafExt{\mathop{\mathcal{E}\!\mathit{xt}}\nolimits}
\def\Sch{\mathit{Sch}}
\def\Mor{\mathop{\mathrm{Mor}}\nolimits}
\def\Ob{\mathop{\mathrm{Ob}}\nolimits}
\def\Sh{\mathop{\mathit{Sh}}\nolimits}
\def\NL{\mathop{N\!L}\nolimits}
\def\CH{\mathop{\mathrm{CH}}\nolimits}
\def\proetale{{pro\text{-}\acute{e}tale}}
\def\etale{{\acute{e}tale}}
\def\QCoh{\mathit{QCoh}}
\def\Ker{\mathop{\mathrm{Ker}}}
\def\Im{\mathop{\mathrm{Im}}}
\def\Coker{\mathop{\mathrm{Coker}}}
\def\Coim{\mathop{\mathrm{Coim}}}

% Boxtimes
%
\DeclareMathSymbol{\boxtimes}{\mathbin}{AMSa}{"02}

%
% Macros for moduli stacks/spaces
%
\def\QCohstack{\mathcal{QC}\!\mathit{oh}}
\def\Cohstack{\mathcal{C}\!\mathit{oh}}
\def\Spacesstack{\mathcal{S}\!\mathit{paces}}
\def\Quotfunctor{\mathrm{Quot}}
\def\Hilbfunctor{\mathrm{Hilb}}
\def\Curvesstack{\mathcal{C}\!\mathit{urves}}
\def\Polarizedstack{\mathcal{P}\!\mathit{olarized}}
\def\Complexesstack{\mathcal{C}\!\mathit{omplexes}}
% \Pic is the operator that assigns to X its picard group, usage \Pic(X)
% \Picardstack_{X/B} denotes the Picard stack of X over B
% \Picardfunctor_{X/B} denotes the Picard functor of X over B
\def\Pic{\mathop{\mathrm{Pic}}\nolimits}
\def\Picardstack{\mathcal{P}\!\mathit{ic}}
\def\Picardfunctor{\mathrm{Pic}}
\def\Deformationcategory{\mathcal{D}\!\mathit{ef}}
\usepackage{cleveref}
\crefname{originalreference}{original reference}{original references}
\Crefname{originalreference}{Original reference}{Original references}
\makeatletter
\newlabel{lemma-characterize-finite-locally-constant}{{03RV}{source}{}{sourcepointer-1}{}}
\newlabel{lemma-characterize-finite-locally-constant@cref}{{[originalreference][0][]03RV}{[1][0][]source}{}{}{}}
\makeatother
\crefname{originalreference}{original reference}{original references}
\Crefname{originalreference}{Original reference}{Original references}
\makeatletter
\newlabel{lemma-closed-of-affine-normal-scheme-with-alg-closed-function-field}{{09ZA}{source}{}{sourcepointer-2}{}}
\newlabel{lemma-closed-of-affine-normal-scheme-with-alg-closed-function-field@cref}{{[originalreference][0][]09ZA}{[1][0][]source}{}{}{}}
\makeatother
\crefname{originallemma}{original lemma}{original lemmas}
\Crefname{originallemma}{Original lemma}{Original lemmas}
\makeatletter
\newlabel{sites-cohomology-lemma-torsors-h1}{{03AJ}{source}{}{sourcepointer-3}{}}
\newlabel{sites-cohomology-lemma-torsors-h1@cref}{{[originallemma][0][]03AJ}{[1][0][]source}{}{}{}}
\makeatother
\title{The first etale cohomology of a finite constant group on a closed subscheme of a normal affine scheme with separably closed function field vanishes}
\author{The Stacks Project Authors}
\date{stacks-project:etale-cohomology.tex:local-snapshot-20261011; GNU FDL 1.2 or later}
\begin{document}
\maketitle
\noindent\textit{Editorial note.} All delivered mathematical text is the author's own, copied from the pinned Stacks Project chapter file etale-cohomology.tex; no statement or proof step was rewritten, reordered or removed.\par\medskip
\noindent\textit{Editorial note.} The counted proof is the author's complete argument: the translation of a cohomology class into a torsor and its description as a finite locally constant sheaf, the representation of the torsor by a finite etale scheme, and the closing trivialization through the existence of a section.\par\medskip
\noindent\textit{Editorial note.} The label lemma-gabber-for-h1-absolutely-algebraically-closed names the counted statement, and the proof environment immediately following it in the source is delivered verbatim as its proof.\par\medskip
\noindent\textit{Editorial note.} The counted statement is carried by the author's own lemma environment in the chapter source.\par\medskip
\noindent\textit{Editorial note.} Pre-open gate: python operations/agent-authoring/tools/precard\_gate.py --source etale-cohomology.tex --label lemma-gabber-for-h1-absolutely-algebraically-closed --start 14424 --end 14442 --local-id 3572 returned RESULT=OK (L3572 free, occupancy FREE, public interval clear). Candidate taken from operations/pool-cards/stacks-etale-cohomology-halfB.json.\par\medskip
\noindent\textit{Editorial note.} The statement and proof are copied from the pinned chapter file; the cited inputs are kept as external pointers with their canonical Stacks tags.\par\medskip
\noindent\textit{Allowed standard prerequisites (editorial).}\begin{itemize}
\item The first cohomology with constant finite coefficients classifies torsors. Cohomology on Sites, the lemma identifying the first cohomology group with coefficients in a constant finite group with the set of torsors, and describing them as finite locally constant sheaves. Supplies the translation of a cohomology class into a torsor on the closed subscheme, which is a finite locally constant sheaf.
\item Representability of a finite locally constant sheaf by a finite etale scheme. Same chapter, the lemma representing a finite locally constant sheaf by a finite etale morphism of schemes. Supplies the scheme V finite etale over Z representing the torsor, which is surjective over Z because a torsor is locally trivial.
\item Sections of finite etale covers of a closed subscheme of a normal affine scheme. Same chapter, the lemma producing a section of a finite etale cover of a closed subscheme of a normal affine scheme with absolutely algebraically closed function field. Supplies the section of V -> Z, which makes the torsor trivial and hence the class zero.
\end{itemize}\medskip
\ifdefined\hypertarget\hypertarget{sourcepointer-1}{}\fi
\noindent\textit{Source reference.} Original reference 03RV: sources/stacks-etale-cohomology/pinned-original-source.tex, source lines 9994--9994. Prerequisite pointer invoked at proof line(s) 14437; it names a standard fact of the same chapter that the counted argument uses, and every step of the counted proof built on it is retained verbatim in the delivered body.\par\medskip
\ifdefined\hypertarget\hypertarget{sourcepointer-2}{}\fi
\noindent\textit{Source reference.} Original reference 09ZA: sources/stacks-etale-cohomology/pinned-original-source.tex, source lines 14390--14390. Prerequisite pointer invoked at proof line(s) 14440; it names a standard fact of the same chapter that the counted argument uses, and every step of the counted proof built on it is retained verbatim in the delivered body.\par\medskip
\ifdefined\hypertarget\hypertarget{sourcepointer-3}{}\fi
\noindent\textit{Source reference.} Original lemma 03AJ: sources/stacks-chow/tags.txt, source lines 4312--4312. Prerequisite pointer invoked at proof line(s) 14432; it names a standard fact of a sibling chapter that the counted argument uses, and every step of the counted proof built on it is retained verbatim in the delivered body.\par\medskip
\begin{lemma}
\label{lemma-gabber-for-h1-absolutely-algebraically-closed}
Let $X$ be a normal integral affine scheme with separably closed
function field. Let $Z \subset X$ be a closed subscheme.
For any finite abelian group $M$ we have $H^1_\etale(Z, \underline{M}) = 0$.
\end{lemma}

\begin{proof}
By Cohomology on Sites, Lemma \ref{sites-cohomology-lemma-torsors-h1}
an element of $H^1_\etale(Z, \underline{M})$ corresponds to a
$\underline{M}$-torsor $\mathcal{F}$ on $Z_\etale$.
Such a torsor is clearly a finite locally constant sheaf.
Hence $\mathcal{F}$ is representable by a scheme $V$ finite
\'etale over $Z$, Lemma \ref{lemma-characterize-finite-locally-constant}.
Of course $V \to Z$ is surjective as a torsor is locally trivial.
Since $V \to Z$ has a section by
Lemma \ref{lemma-closed-of-affine-normal-scheme-with-alg-closed-function-field}
we are done.
\end{proof}

\end{document}
